\documentclass[11pt]{article}
\usepackage{explorer}
\newcommand{\ExplorerLanguage}{FR}
\newcommand{\ExplorerNumber}{05}
\newcommand{\ExplorerKind}{Leçon et méthode}
\begin{document}
\ExplorerTitle{Première--Terminale · dénombrement}{Chemins, combinaisons et Pascal}
\ExplorerSection{Avant de commencer}
Produit, factorielle et choix sans ordre.\par
\ExplorerSection{Idées et méthode}
Un chemin de (0,0) à (a,b) n’utilisant que des pas vers la droite ou vers le haut comporte \(a + b\) pas. Choisir les positions des a pas horizontaux donne \(C(a+b,a)\) chemins.\par
Pour \(0 \leq  k \leq  n, C(n,k) = n\)!/[k!\((n-k)\)!]. L’ordre des objets choisis ne compte pas. La symétrie \(C(n,k) = C(n,n-k)\) revient à choisir les objets retenus ou ceux exclus.\par
La relation de Pascal \(C(n,k) = C(n-1,k-1) + C(n-1,k)\) partage les choix selon la présence d’un objet fixé. On peut aussi classer les chemins selon leur dernier pas.\par
\ExplorerSection{Exemple expliqué}
De (0,0) à (3,2), il \(y a C(5,3) = 10\) chemins. Pour passer par (1,1), il faut choisir un des deux chemins jusqu’à ce point, puis un des trois chemins jusqu’à l’arrivée : \(2 \times  3 = 6\) chemins.\par
\begin{center}\begin{tikzpicture}[scale=.85]\draw[gray!45](0,0)grid(3,2);\draw[mint,very thick,->](0,0)--(1,0)--(1,1)--(2,1)--(3,1)--(3,2);\fill[violet](1,1)circle(3pt)node[above left]{$(1,1)$};\node[below]at(0,0){$(0,0)$};\node[above]at(3,2){$(3,2)$};\end{tikzpicture}\end{center}
\ExplorerSection{À explorer}
Dessiner un petit quadrillage et écrire dans chaque sommet le nombre de chemins qui y arrivent. Additionner les valeurs du voisin de gauche et du voisin du bas.\par
\end{document}
